% !TeX root = ../../example-en.tex
\section{Pages and images}
\sectioncoverpage[\campusexamplesectiononeimage]{Page layouts, equations, image fitting and captions.}
\note{Start with the body area and type sizes, then demonstrate equations, code and images. The background is the general theme's original geometric artwork.}
\subsection{Page basics}

\begin{frame}[campusfocus]{Focus page}
  \begin{description}[Header]
    \item[Header] Only the frame title is shown.
    \item[Footer] Page number, author and title remain; navigation is hidden.
    \item[Option] Add \texttt{campusfocus} to the frame options.
  \end{description}
\end{frame}
\note{Use the focus style when less navigation is useful. Ordinary body frames retain the subsection and frame title on two header lines.}

\begin{frame}[campusbackground=\campusexampleimage,
  campusbackgroundopacity=.18]{Centered image background}
  \begin{description}[Opacity]
    \item[Crop] Fill the page proportionally, align centers and crop the excess.
    \item[Opacity] \texttt{campusbackgroundopacity=.18}; 0 is transparent, 1 is opaque.
    \item[Scope] The background applies the focus style to this frame only.
  \end{description}
\end{frame}
\note{Adjust opacity to preserve text contrast. This is decorative imagery, so a content caption is unnecessary.}

\usetikzlibrary{positioning,arrows.meta}
\begin{frame}{Aligned flowchart}
  \begin{figure}
    \centering
    \begin{tikzpicture}[
      node distance=1.7cm,
      stage/.style={draw=maincolor,fill=maincolor!8,rounded corners=2pt,
        minimum width=3cm,minimum height=1.2cm,align=center},
      flow/.style={draw=maincolor,line width=.7pt,-{Latex[length=2.2mm]}}]
      \node[stage] (draft) {Free-form draft};
      \node[stage,right=of draft] (outline) {Detailed outline};
      \node[stage,right=of outline] (slides) {Slides and notes};
      \draw[flow] (draft.east) -- (outline.west);
      \draw[flow] (outline.east) -- (slides.west);
    \end{tikzpicture}
    \caption{From a free-form draft to an approved outline, then slides and notes}
  \end{figure}
\end{frame}
\note{Point out aligned nodes and evenly spaced connections. Approve the detailed outline before producing slides and notes; diagram symmetry should follow its meaning.}

\begin{frame}{Draft, outline and production}
  \begin{description}[Produce]
    \item[Prepare] Fill in \texttt{outline.md}; put sources in \texttt{materials/}.
    \item[Outline] The agent develops sections and subsections for your review.
    \item[Produce] The authoring skill creates slides and notes from the approved outline.
    \item[Deliver] The validation skill checks the PDF and PPTX in \texttt{build/}.
  \end{description}
  \medskip
  Ask the agent: Make a presentation from the outline.
\end{frame}
\note{The authoring entry is campus-beamer-authoring; validation uses campus-beamer-validation. Public examples live in examples/sections/ and examples/english/; private talk chapters live in sections/. The user supplies purpose, audience, duration, a free-form draft and local materials. The agent develops explanations, evidence, transitions and limitations in build/outline-normalized.md for approval. Production selects layouts and allocates content between slides and notes. Layout changes can stay within approved content; new substantive claims return to outline review. STYLE.md controls presentation preferences.}

% Diagnostic instrumentation is local to this frame; the theme is unchanged.
\begingroup
  % Horizontal dimensions use mm; vertical dimensions use the named font's
  % baseline grid (header: tiny/Large/normalsize; footer: scriptsize).
\newcommand{\layoutGuideDistance}[2]{%
  \begingroup
    \pgfextractx{\dimen0}{\pgfpointdiff{\pgfpointanchor{#1}{center}}{\pgfpointanchor{#2}{center}}}%
    \pgfextracty{\dimen2}{\pgfpointdiff{\pgfpointanchor{#1}{center}}{\pgfpointanchor{#2}{center}}}%
    \pgfmathparse{veclen(\the\dimen0,\the\dimen2)/(1mm)}%
    \pgfmathprintnumber[fixed,precision=2]{\pgfmathresult}\,mm%
  \endgroup
}
% Vertical dimensions are expressed in the same grid unit used by the
% navigation and footer rows. Horizontal dimensions continue to use mm.
\newcommand{\layoutGuideBaselineDistance}[3][\campusfooterbaseline]{%
  \begingroup
    \pgfextracty{\dimen2}{\pgfpointdiff{\pgfpointanchor{#2}{center}}{\pgfpointanchor{#3}{center}}}%
    \pgfmathparse{abs(\the\dimen2)/(1mm) / ((\the#1)/(1mm))}%
    \pgfmathprintnumber[fixed,precision=2]{\pgfmathresult}\,\mbox{baselineskip}%
  \endgroup
}
\newcommand{\layoutGuideHorizontalBaselineDistance}[3][\campusheadertextgap]{%
  \begingroup
    \pgfextractx{\dimen0}{\pgfpointdiff{\pgfpointanchor{#2}{center}}{\pgfpointanchor{#3}{center}}}%
    \pgfmathparse{abs(\the\dimen0)/(1mm) / ((\the#1)/(1mm))}%
    \pgfmathprintnumber[fixed,precision=2]{\pgfmathresult}\,\mbox{baselineskip}%
  \endgroup
}
\begin{frame}{Alignment guides}
  % Draw last, above Beamer's headline/footline, on this page only.
  \AddToHookNext{shipout/foreground}{%
    \begin{tikzpicture}[remember picture,overlay]
      \tikzset{layout guide/.style={red,line width=.4pt},
        layout label/.style={text=red,font=\tiny,fill=white,inner sep=.7pt},
        layout dimension/.style={layout guide,<->}}
      % Both transparent assets are square and cropped to the emblem bounds.
      \coordinate (emblem-left) at
        ([xshift=\dimexpr\campuslogoleftskip+\campuslogopadding\relax,
          yshift=\dimexpr-\campuslogoheight+\campuslogopadding+.5\campuslogoimagewidth\relax]
          current page.north west);
      \coordinate (emblem-right) at ([xshift=\campuslogoimagewidth]emblem-left);
      \coordinate (emblem-bottom) at
        ([xshift=\dimexpr\campuslogoleftskip+.5\campuslogowidth\relax,
          yshift=\dimexpr-\campuslogoheight+\campuslogopadding\relax]
          current page.north west);
      \draw[layout guide]
        ([yshift=.5\campuslogoimagewidth]emblem-left) rectangle
        ([yshift=-.5\campuslogoimagewidth]emblem-right);
      \draw[layout dimension] ([xshift=-\campuslogopadding]emblem-left) -- (emblem-left);
      \draw[layout dimension] (emblem-right) -- ([xshift=\campuslogopadding]emblem-right);
      \draw[layout dimension] (emblem-bottom) -- ([yshift=-\campuslogopadding]emblem-bottom);
      \node[layout label,anchor=north west] at
        ([xshift=\campuslogoleftskip,yshift=\dimexpr-\campuslogoheight-8mm\relax]
          current page.north west)
        {Emblem padding: \the\campuslogopadding};
      \foreach \offset/\tag in {
        \campuslogoleftskip/V1,
        {\dimexpr\campuslogoleftskip+\campuslogowidth\relax}/V2,
        {\dimexpr\campuslogoleftskip+\campusframetitleleftskip\relax}/V3,
        {\dimexpr\paperwidth-\campuslogoleftskip\relax}/V4}{%
        \coordinate (\tag) at ([xshift=\offset]current page.north west);
        \draw[layout guide] (\tag) -- (\tag |- current page.south);
        \node[layout label,anchor=north] at
          ([yshift=-16mm]\tag) {\tag};
      }
      \draw[layout guide,dashed] (current page.north) -- (current page.south);
      \node[layout label,anchor=north] at ([yshift=-16mm]current page.north) {Center};

      % H1, H2 and H3 are the actual boundaries of the three header grid
      % rows: after the empty row, after the title row, and after the
      % subtitle row.
      \coordinate (H1) at ([yshift=-\campusheaderemptybaseline]current page.north west);
      \coordinate (H2) at ([yshift=-\dimexpr\campusheaderemptybaseline+\campusheaderbaseline\relax]current page.north west);
      \coordinate (H3) at ([yshift=-\dimexpr\campusheaderemptybaseline+\campusheaderbaseline+\campusheadersubtitlebaseline\relax]current page.north west);
      \coordinate (H5) at ([yshift=\dimexpr\campusnavigationfooterheight+\campusfootlineheight+\campusfootlinedepth\relax]current page.south west);
      \coordinate (H6) at ([yshift=\dimexpr\campussectionnavigationheight+\campusfootlineheight+2\campusfootlinedepth\relax]current page.south west);
      \coordinate (H7) at ([yshift=\dimexpr\campusfootlineheight+\campusfootlinedepth\relax]current page.south west);
      \coordinate (H8) at (current page.south west);
      % H4 is the body safety boundary: keep a baseline-grid breathing room
      % above the first footer band (H5).
      \coordinate (H4) at ([yshift=2.4\campusfooterbaseline]H5);
      \foreach \tag in {H1,H2,H3,H4,H5,H6,H7}{%
        \draw[layout guide] (\tag) -- (current page.east |- \tag);
        \node[layout label,anchor=east] at
          ([xshift=-2mm]current page.east |- \tag) {\tag};
      }
      \coordinate (safe-a) at ([xshift=15mm]current page.west |- H4);
      \coordinate (safe-b) at (safe-a |- H5);
      \draw[layout dimension] (safe-a) -- (safe-b)
        node[midway,right,layout label] {\layoutGuideBaselineDistance{safe-a}{safe-b}};
      % An inset half-stroke keeps the bottom boundary visible in the PDF.
      \draw[layout guide] ([yshift=.2pt]current page.south west)
        -- ([yshift=.2pt]current page.south east);

      % Show all three header grid rows separately.  These coordinates are
      % layout boundaries, so glyph ascenders and descenders do not change the
      % reported one-baseline lengths.
      \coordinate (header-top) at (current page.north west);
      \foreach \upper/\lower/\pos/\label/\unit in {
        header-top/H1/.70/Top space (\string\tiny)/\campusheaderemptybaseline,
        H1/H2/.78/Subsection (\string\Large)/\campusheaderbaseline,
        H2/H3/.86/Frame title (\string\normalsize)/\campusheadersubtitlebaseline}{%
        \coordinate (header-a) at ([xshift=\pos\paperwidth]current page.west |- \upper);
        \coordinate (header-b) at (header-a |- \lower);
        \draw[layout dimension] (header-a) -- (header-b)
          node[midway,left,layout label] {\label: \layoutGuideBaselineDistance[\unit]{header-a}{header-b}};
      }

      % Horizontal dimensions occupy the open area below the diagnostic text.
      \coordinate (logo-a) at ([yshift=-63mm]V1);
      \coordinate (logo-b) at (V2 |- logo-a);
      \draw[layout dimension] (logo-a) -- (logo-b)
        node[midway,above,anchor=south west,layout label] {Emblem block: %
          \texttt{1 baselineskip (\string\Huge) + 1 baselineskip (\string\tiny)
          = \the\campuslogowidth}};
      \coordinate (gap-a) at ([yshift=-67mm]V2);
      \coordinate (gap-b) at (V3 |- gap-a);
      \draw[layout dimension] (gap-a) -- (gap-b)
        node[midway,above,anchor=south west,layout label] {%
          \layoutGuideHorizontalBaselineDistance{gap-a}{gap-b} (\string\normalsize)};
      \coordinate (body-a) at ([yshift=-70mm]V1);
      \coordinate (body-b) at (V4 |- body-a);
      \draw[layout dimension] (body-a) -- (body-b)
        node[midway,above,layout label] {Body width: \layoutGuideDistance{body-a}{body-b}};
      \coordinate (margin-a) at ([yshift=-72mm]current page.north west);
      \coordinate (margin-b) at (V1 |- margin-a);
      \draw[layout dimension] (margin-a) -- (margin-b)
        node[midway,above,anchor=south west,layout label] {%
          \layoutGuideHorizontalBaselineDistance[\campusheaderbaseline]{margin-a}{margin-b} (\string\Large)};
      \coordinate (margin-c) at (V4 |- margin-a);
      \coordinate (margin-d) at (current page.east |- margin-a);
      \draw[layout dimension] (margin-c) -- (margin-d)
        node[midway,above,anchor=south east,layout label] {%
          \layoutGuideHorizontalBaselineDistance[\campusheaderbaseline]{margin-c}{margin-d} (\string\Large)};

      % Footer heights are read from the theme, so changes remain in sync.
      % Callouts above the bands keep their text clear of the navigation.
      \foreach \upper/\lower/\pos/\label in {
        H5/H6/.55/subsection (\string\scriptsize),
        H6/H7/.72/section (\string\scriptsize),
        H7/H8/.76/Information (\string\scriptsize)}{%
        \coordinate (band-a) at ([xshift=\pos\paperwidth]current page.west |- \upper);
        \coordinate (band-b) at (band-a |- \lower);
        \draw[layout dimension] (band-a) -- (band-b)
          node[midway,right,layout label] {\label: \layoutGuideBaselineDistance{band-a}{band-b}};
      }
    \end{tikzpicture}%
  }%
  \vspace*{11mm}
  \hspace*{-2pt}%
  \begin{minipage}{\dimexpr\textwidth-6pt\relax}
    \begin{description}[Vertical]
      \item[Vertical] V1: body left; V2: emblem right; V3: header text; V4: body right.
      \item[Header] H1--H3 mark header rows; arrows use each row's baselineskip.
      \item[Footer] H5--H7 mark the footer bands; spacing uses \texttt{\string\scriptsize} baselineskip.
      \item[Guide] H4 marks body clearance above the footer; the dashed line is the center.
    \end{description}
  \end{minipage}
\end{frame}
\note{Use the red guides to inspect emblem, header, body and footer alignment. This diagnostic page is part of the feature tour.}
\endgroup

% Read sizes from the active LaTeX commands rather than duplicating point values.
\begingroup
\makeatletter
\newcommand{\fontGuideSize}[1]{%
  \begingroup #1\edef\fontGuideValue{\f@size}%
    \footnotesize\fontGuideValue\,pt\endgroup}
\makeatother
\newcommand{\fontGuideSample}[1]{%
  \makebox[24mm][l]{{\color{maincolor}#1 Aa}}}
\begin{frame}{Type sizes and roles}
  \footnotesize
  \setlength{\tabcolsep}{0pt}
  \renewcommand{\arraystretch}{1.05}
  \begin{tabular}{@{}l@{\hspace{2mm}}p{28mm}@{\hspace{2mm}}p{15mm}@{\hspace{2mm}}p{72mm}@{}}
    \textbf{Sample} & \textbf{Command} & \textbf{Size} & \textbf{Main roles} \\[.15mm]
    \fontGuideSample{\Huge\bfseries} & \texttt{\string\Huge} & \fontGuideSize{\Huge} &
      Main titles on cover, section and closing pages \\[.15mm]
    \fontGuideSample{\LARGE} & \texttt{\string\LARGE} & \fontGuideSize{\LARGE} &
      Single-line headers; the closing thank-you message \\[.15mm]
    \fontGuideSample{\Large\bfseries} & \texttt{\string\Large} & \fontGuideSize{\Large} &
      Upper body header; section number; cover subtitle \\[.15mm]
    \fontGuideSample{\large} & \texttt{\string\large} & \fontGuideSize{\large} &
      Optional emphasis in body text \\[.15mm]
    \fontGuideSample{\normalsize} & \texttt{\string\normalsize} & \fontGuideSize{\normalsize} &
      Body text; lower header; metadata; subsection guide \\[.15mm]
    \fontGuideSample{\small} & \texttt{\string\small} & \fontGuideSize{\small} &
      Callout body; nested lists; this demo's bibliography \\[.15mm]
    \fontGuideSample{\footnotesize} & \texttt{\string\footnotesize} & \fontGuideSize{\footnotesize} &
      Callout titles; third-level lists; captions \\[.15mm]
    \fontGuideSample{\scriptsize} & \texttt{\string\scriptsize} & \fontGuideSize{\scriptsize} &
      Navigation and the footer information row \\[.15mm]
    \fontGuideSample{\tiny} & \texttt{\string\tiny} & \fontGuideSize{\tiny} &
      Diagnostic labels; corner citations; paper metadata
  \end{tabular}
  \par\smallskip
  {\scriptsize Values are read from this document in TeX points. Weight follows the element style.\par
  Bibliography commands accept a font size; frame shrinking can reduce the final display further.\par}
\end{frame}
\note{Compare the actual sizes. Normal size is the body default, while code uses small. Split crowded content before reducing type sizes.}
\endgroup


\begin{frame}{Three levels of structure}
  \begin{description}[Subsection]
    \item[Section] Groups a chapter and appears on section pages and in navigation.
    \item[Subsection] Groups a topic and supplies the upper header line.
    \item[Frametitle] Describes this page on the smaller, lower header line.
    \item[More pages] Reuse a subsection and give each frame its own title.
  \end{description}
\end{frame}
\note{Distinguish chapter structure, topic labels and individual page titles before moving to specific content examples.}

\begin{frame}{Inline math and derivations}
  For $\mathbf{x}_i\in\mathbb{R}^{d}$, predict the scalar target $y_i$ with a linear model:
  \begin{align*}
    \widehat{y}_i &= \mathbf{w}^{\mathsf T}\mathbf{x}_i+b, \\
    \mathcal{L}(\mathbf{w},b)
      &= \frac{1}{N}\sum_{i=1}^{N}(\widehat{y}_i-y_i)^2
         +\lambda\lVert\mathbf{w}\rVert_2^2, \\
    \nabla_{\mathbf{w}}\mathcal{L}
      &= \frac{2}{N}\sum_{i=1}^{N}(\widehat{y}_i-y_i)\mathbf{x}_i
         +2\lambda\mathbf{w}.
  \end{align*}
  \begin{notebox}[Typesetting]
    Use \texttt{\$...\$} for inline math and \texttt{align*} for multiple lines.
    Put \texttt{\&} before the equals sign; the star suppresses numbering.
  \end{notebox}
\end{frame}
\note{Follow the equals signs through prediction, loss and gradient. The example illustrates aligned mathematical typesetting.}

\begin{frame}{Matrices, cases and notation}
  \begin{columns}[T,onlytextwidth]
    \begin{column}{.47\textwidth}
      \textbf{Matrix: batch prediction}
      \[
        \mathbf{X}=\begin{bmatrix}
          \mathbf{x}_1^{\mathsf T}\\ \vdots\\ \mathbf{x}_N^{\mathsf T}
        \end{bmatrix},\qquad
        \widehat{\mathbf{y}}=\mathbf{X}\mathbf{w}+b\mathbf{1}.
      \]
      \texttt{bmatrix} supplies square brackets; \texttt{\&} separates columns.
    \end{column}
    \begin{column}{.47\textwidth}
      \textbf{Piecewise function: ReLU}
      \[
        f(z)=\begin{cases}z, & z\geq 0,\\ 0, & z<0.\end{cases}
      \]
      \texttt{cases} aligns values and conditions. Use \texttt{\string\text\{...\}} for words inside math.
    \end{column}
  \end{columns}
  \medskip
  \begin{tipbox}[Define the symbols]
    $N$ is the sample count, $d$ the feature dimension, and $\mathbf{1}\in\mathbb{R}^{N}$ the all-ones vector.
    Keep notation and its explanation on the same page.
  \end{tipbox}
\end{frame}
\note{Compare the matrix spacing with the piecewise conditions. Define the symbols where they are first used.}

\begin{frame}[fragile]{Python code}
\begin{campuscode}[language=Python,numbers,caption={Indentation makes the function structure clear.}]{normalize.py}
def normalize(values):
    # Normalize values while preserving clear indentation.
    total = sum(values)
    if total == 0:
        return [0.0 for value in values]
    return [value / total for value in values]

print(normalize([2, 3, 5]))
\end{campuscode}
\small The title bar identifies the file and language; captions use the figure style.
\end{frame}
\note{Point out the Python icon, filename, indentation, syntax highlighting and optional line numbers. This example assumes nonnegative values and returns zeros when their sum is zero. Inline code requires a fragile frame. The optional caption has its own line below the code window and does not increment the figure counter.}

\begin{frame}[fragile]{Terminal commands and output}
\begin{campuscode}[language=bash,caption={A static transcript; commands are not executed during compilation.}]{Terminal}
$ python -m venv .venv
$ source .venv/bin/activate
$ python normalize.py
[0.2, 0.3, 0.5]
\end{campuscode}
\small The prompt distinguishes input from output.
\end{frame}
\note{Lines beginning with a dollar sign are commands; the final line is their displayed output. The activation command targets a typical Unix shell. The icon and label identify the content type and can be overridden separately from syntax highlighting.}

\begin{frame}[fragile]{Read code from a file}
\campusinputcode[language=Python,description={Read and highlight an existing source file.}]
  {normalize.py}{examples/sections/code/normalize.py}
\small Reuse source with \verb|\campusinputcode[language=Python]{name}{path}|.
\end{frame}
\note{This page reads the shared normalize.py source. Use firstline and lastline for a snippet and firstnumber for original line numbers. The description option is an alias for caption.}

\begin{frame}[fragile]{Fit an image}
  \begin{columns}[T,onlytextwidth]
    \begin{column}{.48\textwidth}
      \textbf{Fit and caption}
      \begin{verbatim}
\begin{figure}
  \fitgraphic[
    width=.42,
    height=.40
  ]{\schoolwordmarkonlight}
  \caption{Fitted wordmark}
\end{figure}
      \end{verbatim}
      Width and height are fractions of the page; the aspect ratio is preserved.
    \end{column}
    \begin{column}{.48\textwidth}
      \begin{figure}
        \fitgraphic[width=.42,height=.40]{\schoolwordmarkonlight}
        \caption{Fitted wordmark}
      \end{figure}
    \end{column}
  \end{columns}
\end{frame}
\note{Compare the source on the left with its output on the right. The image stays within both page-relative bounds without distortion.}

\begin{frame}[fragile]{Fit an image with a caption}
  \begin{columns}[T,onlytextwidth]
    \begin{column}{.48\textwidth}
      \textbf{Image and caption together}
      \begin{verbatim}
\fitfigure[
  width=.42,
  height=.28
]{\schoolwordmarkonlight}
 {Fitted horizontal wordmark}
      \end{verbatim}
      Leave the final argument empty to omit the caption.
    \end{column}
    \begin{column}{.48\textwidth}
      \fitfigure[width=.42,height=.28]{\schoolwordmarkonlight}{Fitted horizontal wordmark}
    \end{column}
  \end{columns}
\end{frame}
\note{The fitfigure command combines fitting and caption placement. An empty caption argument leaves no caption line.}

\begin{rightimageframe}[image width=.40]{Full-height image}{\campusexampleimage}
  \begin{description}[Height]
    \item[Height] The image reaches from the page top to the navigation area.
    \item[Width] \texttt{image width} sets its fraction of the page width.
    \item[Body] Lists, equations and callouts remain on the left.
  \end{description}
\end{rightimageframe}
\note{The right-hand image ends above navigation; body content stays in the left column. The talk selects the image explicitly.}

\begin{frame}[fragile]{Image command reference}
  \begin{description}[PDF page]
    \item[Fit] \verb|\fitgraphic[width=.8,height=.6]{file}|
    \item[Caption] \verb|\fitfigure[width=.8,height=.6]{file}{caption}|
    \item[Right] \verb|\begin{rightimageframe}{title}{file}|
    \item[PDF page] Use \texttt{page=2} to select the second page of a PDF.
  \end{description}
\end{frame}
\note{Recap the three image interfaces and PDF page selection before introducing paper citations.}
